
\documentclass{article}
\usepackage{colortbl}
\usepackage{makecell}
\usepackage{multirow}
\usepackage{supertabular}

\begin{document}

\newcounter{utterance}

\centering \large Interaction Transcript for game `codenames', experiment `association\_easy', episode 9 with Qwen3.6{-}35B{-}A3B{-}FP8{-}without{-}reasoning.
\vspace{24pt}

{ \footnotesize  \setcounter{utterance}{1}
\setlength{\tabcolsep}{0pt}
\begin{supertabular}{c@{$\;$}|p{.15\linewidth}@{}p{.15\linewidth}p{.15\linewidth}p{.15\linewidth}p{.15\linewidth}p{.15\linewidth}}
    \# & $\;$A & \multicolumn{4}{c}{Game Master} & $\;\:$B\\
    \hline

    \theutterance \stepcounter{utterance}  
    & & & \multicolumn{2}{p{0.3\linewidth}}{
        \cellcolor[rgb]{0.9,0.9,0.9}{
            \makecell[{{p{\linewidth}}}]{
                \texttt{\tiny{[GM$|$GM]}}
                {'hidden': {'team': ['κόρνο', 'δοκάρι', 'κύμβαλο', 'τοίχος', 'είσοδος', 'χτυπητήρι', 'ταχυδρόμος', 'νοσοκόμα', 'κουτάλα'], 'innocent': ['έρημος', 'σφυρίδα', 'βενζίνη', 'γλυνός', 'λαυράκι', 'πετρέλαιο', 'νησί'], 'opponent': ['ρούμι', 'άλογο', 'τζαζ', 'σαμπάνια', 'καναπές', 'καμηλοπάρδαλη', 'oldies', 'καρέκλα'], 'assassin': ['φλανέλα']}, 'revealed': {'team': {'team': [], 'innocent': [], 'opponent': [], 'assassin': []}, 'opponent': {'team': [], 'innocent': [], 'opponent': [], 'assassin': []}}}
            }
        }
    }
    & & \\ \\

    \theutterance \stepcounter{utterance}  
    & & \multicolumn{4}{p{0.6\linewidth}}{
        \cellcolor[rgb]{0.9,0.9,0.9}{
            \makecell[{{p{\linewidth}}}]{
                \texttt{\tiny{[P1$\langle$GM]}}
                \texttt{Δώστε μία μονολεκτική ένδειξη που σχετίζεται με κάποιες από τις λέξεις της ομάδας σας από τον παρακάτω κατάλογο. Μπορείτε να επιλέξετε σε πόσες και σε ποιες λέξεις αναφέρεται η ένδειξή σας.} \\
\texttt{Ο στόχος είναι να μαντέψει ο συνεργάτης σας όλες τις λέξεις της ομάδας σας με βάση τις ενδείξεις που του δίνετε, πριν η άλλη ομάδα αποκαλύψει όλες τις δικές της. Αν ο συνεργάτης σας μαντέψει μια λέξη του 'αντιπάλου' ή μια λέξη 'αντιπερισπασμού', η σειρά μαντέματός του τελειώνει και η αντίπαλη ομάδα παίρνει τη σειρά της. Αν ο συνεργάτης σας μαντέψει μια λέξη 'θανάτου', το παιχνίδι τελειώνει και η ομάδα σας χάνει αμέσως.} \\
\texttt{Όσο περισσότερες λέξεις σχετίζονται με την ένδειξή σας, τόσο πιο γρήγορα μπορείτε να κερδίσετε το παιχνίδι, αλλά τόσο πιο δύσκολοι μπορεί να είναι οι συσχετισμοί για να τους μαντέψει ο συνεργάτης σας, οπότε επιλέξτε τη λέξη{-}ένδειξή σας με σύνεση. Η λέξη{-}ένδειξη πρέπει να σχετίζεται σημασιολογικά με τις λέξεις{-}στόχους, δεν μπορεί να είναι μία από τις λέξεις των καταλόγων ή να περιέχει μέρη τους.} \\
\\ 
\texttt{Να δίνετε πάντα τη μονολεκτική σας ένδειξη και τον κατάλογο των σχετικών λέξεων{-}στόχων χωρισμένο με κόμματα με την ακόλουθη μορφή και να κάνετε τις απαντήσεις σας όσο το δυνατόν πιο σύντομες, μην περιλαμβάνετε ποτέ άλλο κείμενο από αυτό που απαιτείται σε αυτή τη μορφή:} \\
\\ 
\texttt{ΕΝΔΕΙΞΗ: <WORD>} \\
\texttt{ΣΤΟΧΟΙ: <TARGETS>} \\
\\ 
\texttt{Οι λέξεις της ομάδας σας είναι: κόρνο, δοκάρι, κύμβαλο, τοίχος, είσοδος, χτυπητήρι, ταχυδρόμος, νοσοκόμα, κουτάλα.} \\
\texttt{Οι λέξεις της ομάδας του αντιπάλου σας είναι: ρούμι, άλογο, τζαζ, σαμπάνια, καναπές, καμηλοπάρδαλη, oldies, καρέκλα.} \\
\texttt{Οι λέξεις αντιπερισπασμού είναι: έρημος, σφυρίδα, βενζίνη, γλυνός, λαυράκι, πετρέλαιο, νησί.} \\
\texttt{Οι λέξεις θανάτου είναι: φλανέλα.} \\
            }
        }
    }
    & \\ \\

    \theutterance \stepcounter{utterance}  
    & \multicolumn{4}{p{0.6\linewidth}}{
        \cellcolor[rgb]{0.9,0.9,0.9}{
            \makecell[{{p{\linewidth}}}]{
                \texttt{\tiny{[P1$\rangle$GM]}}
                \texttt{ΕΝΔΕΙΞΗ: μουσικό} \\
\texttt{ΣΤΟΧΟΙ: κόρνο, κύμβαλο, χτυπητήρι} \\
            }
        }
    }
    & & \\ \\

    \theutterance \stepcounter{utterance}  
    & & & \multicolumn{2}{p{0.3\linewidth}}{
        \cellcolor[rgb]{0.9,0.9,0.9}{
            \makecell[{{p{\linewidth}}}]{
                \texttt{\tiny{[GM$|$GM]}}
                {'player': 'Cluegiver', 'type': 'prefix error', 'utterance': 'ΕΝΔΕΙΞΗ: μουσικό\nΣΤΟΧΟΙ: κόρνο, κύμβαλο, χτυπητήρι', 'message': 'Your clue did not start with the correct prefix (ΥΠΟΔΕΙΞΗ: ).', 'prefix': 'ΥΠΟΔΕΙΞΗ: '}
            }
        }
    }
    & & \\ \\

\end{supertabular}
}

\end{document}
